\begin{chapintro}
系統的レビューでは、問いに合う試験を、偏りなく、できるだけ漏れなく集めます。この章では、データベースの選び方、検索式の組み立て方、2 人で抄録を選別する手順を説明し、架空の抄録を使って組み入れの判断を練習します。最後に、選別の経過を PRISMA のフロー図にまとめます。
\end{chapintro}

\section{どこを探すか}
\label{sec:03-1}

一つのデータベースだけを検索すると、そこに収録されていない雑誌の試験が抜け落ちます。介入のレビューでは、少なくとも次の 3 つを検索するのが標準です。

\begin{description}
\item[MEDLINE（PubMed）] 米国国立医学図書館のデータベースで、無料で使えます。医学の主要な雑誌はほぼ収録されています。
\item[Embase] 欧州の雑誌や、薬剤に関する雑誌の収録が厚い有料のデータベースです。MEDLINE と重複しない論文がかなり含まれます。
\item[CENTRAL] Cochrane が集めた比較試験の登録簿で、他のデータベースから抽出した試験や、学会抄録の試験も含まれます。
\end{description}

日本語の論文を探すなら医中誌 Web を加えます。さらに、試験登録のデータベース（ClinicalTrials.gov、日本の jRCT、WHO の ICTRP）も検索します。試験登録には、計画段階で試験が登録されているので、結果が論文にならなかった試験も見つかります。\kkref{sec:01-4} で述べた出版バイアスを小さくするには、この試験登録の検索が役に立ちます。学会抄録や報告書のように、雑誌に載らない文献（灰色文献）も、必要に応じて探します。

\section{検索式を組み立てる}
\label{sec:03-2}

検索式は、PICO の要素ごとに言葉を集めて作ります。一つの要素の中では同義語を OR でつなぎ、要素どうしは AND でつなぎます。OR は「どれか一つでも含む」、AND は「すべてを含む」という意味なので、要素の中で候補を広げ、要素の間で絞り込むことになります。

言葉には 2 種類を使います。一つは、データベースが論文に付けている統制語です。MEDLINE では MeSH（Medical Subject Headings）と呼ばれ、たとえば Hypertension という MeSH を指定すると、論文の中で high blood pressure と書かれていても拾えます。もう一つは、題名や抄録に出てくる自由語です。新しい論文はまだ統制語が付いていないことがあるので、両方を OR でつなぎます。

アウトカム（脳卒中）は、ふつう検索式に入れません。抄録には主要アウトカムしか書かれていないことが多く、脳卒中を副次アウトカムとして報告した試験が漏れてしまうからです。比較（プラセボ）も同じ理由で入れません。検索式は、図書館の司書や検索の専門家に見てもらうと、漏れが大きく減ります。

\subsection{\hatten PubMed の検索式の書き方}
\label{sec:03-2-1}

降圧薬 A は架空の薬剤なので、形だけの例を示します。PubMed では次のように書きます。

\begin{rcode}
#1 "drug A"[tiab] OR "XY-123"[tiab]
#2 "Hypertension"[MeSH] OR hypertens*[tiab]
   OR "high blood pressure"[tiab]
#3 "randomized controlled trial"[pt]
   OR randomi*[tiab] OR placebo[tiab]
#4 #1 AND #2 AND #3
\end{rcode}

[tiab] は題名と抄録から探す指定、[pt] は論文の種類の指定、末尾の * は語尾が何でもよいという指定（hypertension も hypertensive も拾う）です。\#3 は RCT に絞るための組み合わせで、Cochrane が漏れの少ない検索式（高感度検索式）を公開しています。この RCT フィルターは MEDLINE のように試験以外の論文も大量に含むデータベースに使うもので、CENTRAL には使いません（Cochrane ハンドブック 4.4.7）。CENTRAL は比較試験を集めた登録簿で、フィルターをかけると本来拾える試験を落とす恐れがあるからです。

\section{重複を除き、2 人で選別する}
\label{sec:03-3}

複数のデータベースを検索すると、同じ論文が何度も出てきます。文献管理ソフトで重複を除いてから選別に入ります。

選別は 2 段階で行います。最初に題名と抄録だけを読み、明らかに基準に合わないものを除きます。残ったものについて本文を入手し、基準に合うかを確かめます。Cochrane ハンドブック（4.6 と MECIR の基準）は、抄録の段階の選別は 1 人でもよいとし、本文による最終的な適格性の判定は 2 人が独立に行うよう求めています。本書は、どちらの段階でも 2 人がそれぞれ独立に判断することを勧めます。判断が分かれたものは話し合うか、3 人目の判断で決めます。1 人で選別すると、判断の誤りや思い込みに気づく機会がなくなるからです。

\begin{hattenbox}
2 人の判断がどれだけ一致したかは κ（カッパ）係数という指標で示すことがあります。偶然の一致を差し引いた一致の度合いで、1 に近いほどよく一致していることを表します。
\end{hattenbox}

抄録の段階では、迷ったものは残します。抄録で除いた文献はふつう二度と見ないので、ここでの除外は本文の段階より慎重に行います。選別の作業には Rayyan のような専用のツールが使えます。

\begin{hattenbox}
選別を機械学習で助ける方法もあります。関連が高いと推定した記録を先に並べ替えて、読む順序を工夫する使い方です。一方、機械が除外と判断した記録を人が見直さないままにすることは、Cochrane ハンドブック（4.6.6）では勧められていません。自動化の道具を使ったときは、その方法をフロー図とあわせて方法の節に書きます。
\end{hattenbox}

\section{抄録の選別を練習する}
\label{sec:03-4}

検索で見つかった架空の抄録を 6 件示します。\kkref{sec:02-3} で決めた基準（高血圧の成人、降圧薬 A とプラセボを比べた RCT、脳卒中を報告、追跡 1 年以上）に照らして、組み入れるか、除外するか、本文で確かめるかを考えてから、続く解説を読んでください。

\begin{description}\small
\item[抄録 1] Garcia 2018（架空）。65 歳以上の高血圧患者 3,200 人を降圧薬 A 群とプラセボ群に無作為に割り付け、中央値 4 年追跡した。主要評価項目の脳卒中は、降圧薬 A 群で有意に少なかった。
\item[抄録 2] Kimura 2017（架空）。診療記録のデータベースから、降圧薬 A を処方された 12,000 人と処方されなかった 30,000 人を抽出し、脳卒中の発症を比べた。降圧薬 A の処方は脳卒中の低いリスクと関連していた。
\item[抄録 3] Lee 2016（架空）。高血圧患者 800 人を、降圧薬 A 群とアムロジピン群に無作為に割り付けた。3 年間の脳卒中の発症に差はなかった。
\item[抄録 4] Novak 2013（架空）。高血圧患者 240 人を降圧薬 A 群とプラセボ群に無作為に割り付け、24 週後の収縮期血圧を比べた。降圧薬 A 群で 12 mmHg 低かった。
\item[抄録 5] Garcia 2020（架空）。Garcia 2018 の試験の参加者を対象に、5 年目までの追跡結果と腎機能の変化を報告する。
\item[抄録 6] Endo 2015（架空）。75 歳以上の高血圧患者 180 人を降圧薬 A 群とプラセボ群に無作為に割り付けた。主要評価項目は心血管イベントで、降圧薬 A 群で少ない傾向がみられた。
\end{description}

抄録 1 は、基準をすべて満たしているので組み入れます。抄録 2 は、診療記録から薬剤を飲んだ人と飲まなかった人を比べた観察研究です。処方されるかどうかは医師が患者の状態を見て決めるので、両群の背景が違い、その違いが結果に混ざります。基準どおり除外します。抄録 3 は RCT ですが、比較が他の降圧薬なので除外します。抄録 4 は、アウトカムが血圧だけで追跡も 24 週と短いので除外します。

抄録 5 は判断に迷う例です。新しい試験ではなく、抄録 1 と同じ試験の追加報告です。同じ参加者を 2 回数えると、その試験の重みが二重になるので、別の試験として組み入れてはいけません。一つの試験に複数の論文がある場合は、試験を単位にしてまとめ、データはどの論文から取るかを決めます。この例では、主要な結果は抄録 1 の論文から取り、5 年目の結果が必要なら抄録 5 の論文を参照します。

抄録 6 は、主要評価項目が「心血管イベント」で、脳卒中を別に報告しているかが抄録からは分かりません。こうした文献は抄録の段階では除かず、本文を入手して確かめます。本文に脳卒中の数が示されていれば組み入れ、なければ著者に問い合わせます。

\section{選別の経過をフロー図にまとめる}
\label{sec:03-5}

選別の経過は、何件を探し、どの段階で何件を、どんな理由で除いたかを、PRISMA 2020（Page 2021）の形式のフロー図にまとめます。降圧薬 A のレビューの架空の例を図 3.1 に示します。

\begin{kkcenter}
\begin{kkfitpic}
\begin{tikzpicture}[
  fb/.style={kplain, minimum width=50mm, text width=46mm, align=left, font=\scriptsize, minimum height=7mm},
  xb/.style={kbox, minimum width=50mm, text width=46mm, align=left, font=\scriptsize, minimum height=7mm},
  node distance=4mm]
\node[fb] (id) {データベース 1,180 件\\試験登録 68 件};
\node[xb, right=6mm of id] (dup) {重複として除外 346 件};
\node[fb, below=of id] (sc) {題名と抄録で選別 902 件};
\node[xb, right=6mm of sc] (ex1) {除外 851 件};
\node[fb, below=of sc] (rt) {本文の入手を試みた 51 報告};
\node[xb, right=6mm of rt] (nr) {本文を入手できず 2 報告};
\node[fb, below=12mm of rt] (el) {本文で適格性を確認 49 報告};
\node[xb, anchor=north west] (ex2) at ([xshift=6mm]el.north east) {除外 33 報告\\対照が実薬 12 報告\\観察研究 9 報告\\脳卒中の報告なし 8 報告\\追跡 1 年未満 4 報告};
\node[fb, below=16mm of el, fill=kkpale] (in) {組み入れた試験 10\\（報告 16 件）};
\foreach \a/\b in {id/sc,sc/rt,rt/el,el/in} \draw[kar] (\a) -- (\b);
\foreach \a/\b in {id/dup,sc/ex1,rt/nr} \draw[kar] (\a) -- (\b);
\draw[kar] (el.east) -- (el.east -| ex2.west);
\end{tikzpicture}
\end{kkfitpic}
\kkcap{図 3.1　降圧薬 A のレビューの PRISMA 2020 フロー図（架空の数）}
\end{kkcenter}

図 3.1 では、本文の段階から数える単位を、文献（報告）の数にしています。同じ試験について複数の報告が出ていることがあり、その場合は試験の数と報告の数が一致しないからです。本文を確認した 49 報告のうち、除外したのは 33 報告（対照が実薬 12、観察研究 9、脳卒中の報告なし 8、追跡 1 年未満 4）です。残りの 16 報告は 10 の試験についての報告で、1 試験につき主要な報告が 1 本、同じ試験の追加報告が合わせて 6 本あります（49 − 33 = 16、10 + 6 = 16）。PRISMA 2020 は、組み入れた試験の数と、その報告の数を別々に示すよう求めています。同じ試験の追加報告は、除外した文献ではなく、組み入れた試験にまとめた報告として数えます。

フロー図を見れば、読者は検索がどこまで及んだか、どんな理由の除外が多かったかを一目で確かめられます。本文の段階で除外した文献は、理由とともに一覧表にして付録に載せるのがふつうです。抄録の段階の除外は件数だけで構いません。

\section{この章のまとめ}
\label{sec:03-6}

検索は MEDLINE、Embase、CENTRAL と試験登録を基本にし、PICO の要素ごとに統制語と自由語を OR でつなぎ、要素どうしを AND でつなぎます。アウトカムは検索式に入れません。選別は抄録と本文の 2 段階で、2 人が独立に判断し、同じ試験の重複報告は試験単位にまとめます。経過は PRISMA 2020 のフロー図に残します。次の章では、組み入れた 10 試験の論文から数字を取り出し、それぞれの試験の偏りを評価します。
